\documentclass{report}
\usepackage{amsfonts, amsmath}

\newcommand\N{{\mathbb N}}
\newcommand\Q{{\mathbb Q}}
\newcommand\R{{\mathbb R}}
\newcommand\Z{{\mathbb Z}}

\pagestyle{empty}
\begin{document}
\large
\centerline{\Huge \bf Analisi Matematica II}
\vskip.2cm
\centerline{\Huge  Quarto Appello (28-01-2003)}
\renewcommand{\thefootnote}{ } 
\footnote{\sl Avete 3 ore di tempo. Ogni esercizio
vale sei punti. La sufficienza si ottiene con un punteggio $\geq$
18. {\bf Solo le risposte chiaramente giustificate saranno prese in
considerazione.} {\scshape Elaborati illeggibili o confusi verranno ignorati.}
{\sffamily La
sintesi sar\`a particolarmente apprezzata.}}
\renewcommand{\thefootnote}{\arabic{footnote}} 

\vskip1cm
\begin{enumerate}
\item Si dimostri che, per ogni $k\in\N$,
$0\leq \frac 1k-\int_k^{k+1}\frac 1\xi d\xi\leq \frac 1{k^2}.$

Usando tale fatto si dimostri che  la successione 
\[
a_n:=\sum_{k=1}^{n}\frac 1k-\int_{1}^{n+1}\frac 1\xi d\xi
\]
ammette limite $\alpha:=\lim_{n\to\infty}a_n$. Infine si mostri
che $\alpha\in(0,2)$.


\item  Si calcoli il seguente integrale
\[
\int_{\frac\pi 4}^{\frac\pi 2}\frac 1{\sin x}dx
\]
\item  Si calcoli il seguente integrale
\[
\int_0^1\frac 1{1+e^t}dt. 
\]
\item Sia $f_n(x):=\sum_{k=1}^n\frac 1k \sin kx$. Si
calcoli, per ogni $n,m\in \N$,
\[
\int_0^{2\pi}f_n(x)f_m(x) dx.
\]

\item Si risolva il seguente problema di Chauchy
\[
y'(x)=4x+xy^2\;;\quad
y(0)=0
\]
\item Si calcoli il raggio di convergenza della serie
\[
\sum_{n=0}^\infty \frac {x^n}{(2n)!}.
\]
\end{enumerate}
\end{document}
